\section{应用需求与框架范围}
联网应用通常不仅需要知道服务的位置。无人机操作员可能需要一架能在截止时间前覆盖指定区域的相机，而推理请求者可能要求特定模型、足够的加速器内存及可接受的准备时间。服务通告只能指出潜在执行者，并不表明它当前愿意且能够承担这项任务。NDNSF 区分提供服务的权限、当前参与意愿、一次调用中的执行者选定，以及任务完成。

即使应用通常只使用一台设备，这些区分仍然重要。操作员选定无人机后，希望后续相机命令作用于同一架无人机；改用另一台可用 Provider 会改变命令含义。相反，图像分析请求可以允许任何满足条件的获准计算 Provider 执行。可复用框架需要支持这两种情形，而不要求应用分别实现无关的通信协议。

发现描述包含参与决策所需的信息。感知任务可能需要粗略区域、时间窗口和传感器要求；推理任务可能需要模型身份、预计输入大小和截止时间。完整图像、详细飞行指令和推理提示词可以在选定后交付。这一边界取决于决策需要：判断可行性不可缺少的信息，不能全部推迟到判断之后。应用仍应限制披露精度。例如，获准相机 Provider 可能需要大致搜索区域，而未选中的候选不必获得完整任务路线。描述加密限制获准服务 Provider 之外的读取者，但不能阻止合法接收者获知或主动分享其解密的信息。

NDNSF-UAV 和 NDNSF-DI 提供互补需求。前者包含设备专属操作、动态观测、实时数据和物理动作；后者包含计算选择、模型准备、大对象、中间张量和增量输出。两者都需要授权与有界请求生命周期。这些差异用于检验共同调用和命名对象接口在不同工作负载下是否仍然适用。可复用性应通过共享机制及应用专用代码来评价，而不能仅凭实现两个演示就认定成立。

\begin{table}[htbp]\centering\small
\caption{框架需求及其与研究问题的关系。}
\begin{tabularx}{\textwidth}{p{.23\textwidth}X p{.16\textwidth}}
\toprule 需求类别 & 所需行为 & 研究关联\\\midrule
发现与调用 & 寻找愿意参与的候选，选定执行者，关联完成或超时 & RQ1、RQ3\\\midrule
授权与保密性 & 检查双方服务权限，限制 Content 访问，应用权限变更 & RQ1、RQ3\\\midrule
执行与协作 & 将工作关联到活动请求；扩展中实施角色和依赖约束 & RQ2、RQ3\\\midrule
命名对象交付 & 区分控制消息、持续流与完整对象，验证取得的 Data & RQ2、RQ3\\\bottomrule
\end{tabularx}
\end{table}

这些是相互关联的需求，而不是严格分离的四个协议层。授权影响远程操作，执行状态决定有效对象能否被使用。后续章节依次介绍 NDN 基础、框架机制与应用验证。研究问题则聚焦于可检验的结论：机制选择是否合理、协作扩展提供何种额外约束，以及各项代价是多少。
